\chapter{2010年新课标卷（文）}

\section{单选题}

\begin{problem}
已知集合 \(A = \{x \mid |x| \leqslant 2, x \in \R\}\)，\(B = \{x \mid \sqrt{x} \leqslant 4, x \in \Z\}\)，则 \(A \cap B =\)（\quad）

\choices
    {\((0,2)\)}
    {\([0,2]\)}
    {\(\{0,2\}\)}
    {\(\{0,1,2\}\)}
\end{problem}

\begin{answer}
D
\end{answer}

\begin{solution}
由 \(|x|\leqslant 2\) 得 \(-2\leqslant x\leqslant 2\)，所以 \(A=[-2,2]\). 又因为 \(\sqrt{x}\) 有意义且 \(\sqrt{x}\leqslant4\)，得 \(0\leqslant x\leqslant16\)；结合 \(x\in\Z\)，有 \(B=\{0,1,2,\ldots,16\}\). 因而 \(A\cap B=\{0,1,2\}\)，故选 D.
\end{solution}

\begin{problem}
已知 \(\bs{a}\)，\(\bs{b}\) 为平面向量，若 \( \bs{a} = (4, 3) \)，\( 2\bs{a} + \bs{b} = (3, 18) \)，则 \(\bs{a}\)，\(\bs{b}\) 夹角的余弦值等于（\quad）

\choices
    {\(\frac{8}{65}\)}
    {\(-\frac{8}{65}\)}
    {\(\frac{16}{65}\)}
    {\(-\frac{16}{65}\)}
\end{problem}

\begin{answer}
C
\end{answer}

\begin{solution}
由 \(2\bs{a}+\bs{b}=(3,18)\) 和 \(\bs{a}=(4,3)\)，得
\[
\bs{b}=(3,18)-2(4,3)=(-5,12)
\]
于是 \(\bs{a}\cdot\bs{b}=4\times(-5)+3\times12=16\)，且 \(|\bs{a}|=5\)，\(|\bs{b}|=13\). 因此 \(\bs{a}\)，\(\bs{b}\) 夹角的余弦为
\[
\frac{\bs{a}\cdot\bs{b}}{|\bs{a}||\bs{b}|}=\frac{16}{5\times13}=\frac{16}{65}
\]
故选 C.
\end{solution}

\begin{problem}
已知复数 \(z = \frac{\sqrt{3} + \i}{(1 - \sqrt{3}\i)^2}\)，则 \(|z| =\)（\quad）

\choices
    {\(\frac{1}{4}\)}
    {\(\frac{1}{2}\)}
    {\(1\)}
    {\(2\)}
\end{problem}

\begin{answer}
B
\end{answer}

\begin{solution}
由复数模的运算性质，
\(|\sqrt{3}+\i|=\sqrt{3+1}=2\)，\(|1-\sqrt{3}\i|=\sqrt{1+3}=2\).
所以
\[
|z|=\frac{|\sqrt{3}+\i|}{|1-\sqrt{3}\i|^{2}}=\frac{2}{2^{2}}=\frac{1}{2}
\]
故选 B.
\end{solution}

\begin{problem}
曲线 \( y = x^3 - 2x + 1 \) 在点 \((1,0)\) 处的切线方程为（\quad）

\choices
    {\(y = x - 1\)}
    {\( y = -x + 1 \)}
    {\(y = 2x - 2\)}
    {\(y = -2x + 2\)}
\end{problem}

\begin{answer}
A
\end{answer}

\begin{solution}
设曲线函数为 \(y=f(x)=x^{3}-2x+1\)，则
\(f'(x)=3x^{2}-2\)，\(f'(1)=1\).
所以曲线在 \((1,0)\) 处的切线斜率为 \(1\). 由点斜式，切线方程为
\[
y-0=1(x-1)
\]
即 \(y=x-1\)，故选 A.
\end{solution}

\begin{problem}
中心在原点，焦点在 \(x\) 轴上的双曲线的一条渐近线经过点 \((4,2)\)，则它的离心率为（\quad）

\choices
    {\(\sqrt{6}\)}
    {\(\sqrt{5}\)}
    {\(\frac{\sqrt{6}}{2}\)}
    {\(\frac{\sqrt{5}}{2}\)}
\end{problem}

\begin{answer}
D
\end{answer}

\begin{solution}
设双曲线方程为
\(\frac{x^{2}}{a^{2}}-\frac{y^{2}}{b^{2}}=1\)，\((a>0,b>0)\).
其渐近线方程为 \(y=\pm\dfrac{b}{a}x\). 由于点 \((4,2)\) 在其中一条渐近线上，得
\[
\frac{b}{a}=\frac{2}{4}=\frac{1}{2}
\]
双曲线的半焦距满足 \(c^{2}=a^{2}+b^{2}\)，所以离心率
\[
e=\frac{c}{a}=\sqrt{1+\frac{b^{2}}{a^{2}}}=\sqrt{1+\frac14}=\frac{\sqrt{5}}{2}
\]
故选 D.
\end{solution}

\begin{problem}
如图，质点 \(P\) 在半径为 \(2\) 的圆周上逆时针运动，其初始位置为 \(P_{0}(\sqrt{2}, -\sqrt{2})\)；角速度为 \(1\)，那么点 \(P\) 到 \(x\) 轴的距离 \(d\) 关于时间 \(t\) 的函数图象大致为
\bitmapfigure[max width=0.35\linewidth]{img/2010/new_curriculum_liberal/q6_condition.png}

\choices
    {\choicebitmap[width=0.15\paperwidth]{img/2010/new_curriculum_liberal/q6_optA.png}}
    {\choicebitmap[width=0.15\paperwidth]{img/2010/new_curriculum_liberal/q6_optB.png}}
    {\choicebitmap[width=0.15\paperwidth]{img/2010/new_curriculum_liberal/q6_optC.png}}
    {\choicebitmap[width=0.15\paperwidth]{img/2010/new_curriculum_liberal/q6_optD.png}}
\end{problem}

\begin{answer}
C
\end{answer}

\begin{solution}
质点初始位置的极角为 \(-\dfrac{\pi}{4}\)，角速度为 \(1\)，故时刻 \(t\) 的纵坐标为
\[
y(t)=2\sin\left(t-\frac{\pi}{4}\right)
\]
点到 \(x\) 轴的距离为纵坐标的绝对值，因此
\[
d(t)=2\left|\sin\left(t-\frac{\pi}{4}\right)\right|
\]
特别地，\(d(0)=\sqrt{2}\)，且在 \(0<t<\dfrac{\pi}{4}\) 时距离递减；\(t=\dfrac{\pi}{4}\) 时 \(d=0\)，\(t=\dfrac{3\pi}{4}\) 时 \(d=2\). 该图象与选项 C 相符，故选 C.
\end{solution}

\begin{problem}
设长方体的长，宽，高分别为 \(2a\)，\(a\)，\(a\)，其顶点都在一个球面上，则该球的表面积为（\quad）

\choices
    {\(3\pi a^{2}\)}
    {\(6\pi a^{2}\)}
    {\(12\pi a^{2}\)}
    {\(24\pi a^{2}\)}
\end{problem}

\begin{answer}
B
\end{answer}

\begin{solution}
长方体的体对角线长为
\[
\sqrt{(2a)^{2}+a^{2}+a^{2}}=\sqrt{6}\,a
\]
球心是长方体体对角线的中点，所以球半径为 \(R=\dfrac{\sqrt{6}}{2}a\). 因而球的表面积为
\[
4\pi R^{2}=4\pi\left(\frac{\sqrt{6}}{2}a\right)^{2}=6\pi a^{2}
\]
故选 B.
\end{solution}

\begin{problem}
如果执行如图所示的框图，输入 \(N = 5\)，则输出的数等于（\quad）
\bitmapfigure[width=0.38\linewidth]{img/2010/new_curriculum_liberal/q8_options.png}

\choices
    {\(\frac{5}{3}\)}
    {\(\frac{4}{5}\)}
    {\(\frac{6}{5}\)}
    {\(\frac{5}{6}\)}
\end{problem}

\begin{answer}
D
\end{answer}

\begin{solution}
框图开始时 \(k=1\)，\(S=0\). 每次先把 \(\dfrac{1}{k(k+1)}\) 加入 \(S\)，再判断 \(k<N\)；若判断为“是”，令 \(k=k+1\) 后继续循环. 因为 \(N=5\)，循环中实际取 \(k=1\)，\(2\)，\(3\)，\(4\)，\(5\)，所以
\[
S=\sum_{k=1}^{5}\frac{1}{k(k+1)}
=\sum_{k=1}^{5}\left(\frac{1}{k}-\frac{1}{k+1}\right)
=1-\frac{1}{6}=\frac{5}{6}
\]
故选 D.
\end{solution}

\begin{problem}
设偶函数 \( f(x) \) 满足 \( f(x) = 2^{x} - 4 \)（\(x \geqslant 0\)），则 \( \{x \mid f(x - 2) > 0\} =\)（\quad）

\choices
    {\(\{x\mid x < -2 \text{或} x > 4\}\)}
    {\(\{x\mid x < 0 \text{或} x > 4\}\)}
    {\(\{x\mid x < 0 \text{或} x > 6\}\)}
    {\(\{x\mid x < -2 \text{或} x > 2\}\)}
\end{problem}

\begin{answer}
B
\end{answer}

\begin{solution}
由于 \(f\) 是偶函数，令 \(u=|x-2|\geqslant0\)，则
\[
f(x-2)=f(|x-2|)=2^{|x-2|}-4
\]
因此
\[
f(x-2)>0
\Longleftrightarrow 2^{|x-2|}>4=2^{2}
\Longleftrightarrow |x-2|>2
\]
解得 \(x<0\) 或 \(x>4\)，所以解集为 \(\{x\mid x<0\text{或}x>4\}\)，故选 B.
\end{solution}

\begin{problem}
若 \(\cos \alpha = -\frac{4}{5}\)，\(\alpha\) 是第三象限的角，则 \(\sin \left( \alpha + \frac{\pi}{4} \right) =\)（\quad）

\choices
    {\(-\frac{7\sqrt{2}}{10}\)}
    {\(\frac{7\sqrt{2}}{10}\)}
    {\(-\frac{\sqrt{2}}{10}\)}
    {\(\frac{\sqrt{2}}{10}\)}
\end{problem}

\begin{answer}
A
\end{answer}

\begin{solution}
因为 \(\alpha\) 是第三象限的角，故 \(\sin\alpha<0\). 由
\[
\sin^{2}\alpha=1-\cos^{2}\alpha=1-\frac{16}{25}=\frac{9}{25}
\]
得 \(\sin\alpha=-\dfrac{3}{5}\). 于是
\[
\sin\left(\alpha+\frac{\pi}{4}\right)
=\sin\alpha\cos\frac{\pi}{4}+\cos\alpha\sin\frac{\pi}{4}
=\left(-\frac35-\frac45\right)\frac{\sqrt{2}}{2}
=-\frac{7\sqrt{2}}{10}
\]
故选 A.
\end{solution}

\begin{problem}
已知平行四边形 \(ABCD\) 的三个顶点为 \(A(-1,2)\)，\(B(3,4)\)，\(C(4,-2)\)，点 \((x,y)\) 在平行四边形 \(ABCD\) 的内部，则 \(z = 2x - 5y\) 的取值范围是（\quad）

\choices
    {\((-14, 16)\)}
    {\((-14, 20)\)}
    {\((-12, 18)\)}
    {\((-12, 20)\)}
\end{problem}

\begin{answer}
B
\end{answer}

\begin{solution}
由平行四边形对角线关系，第四个顶点为
\[
D=A+C-B=(-1,2)+(4,-2)-(3,4)=(0,-4)
\]
令 \(z=2x-5y\)，在四个顶点处分别有
\(z_A=-12\)，\(z_B=-14\)，\(z_C=18\)，\(z_D=20\).
线性函数在平行四边形上的最大值和最小值分别在顶点处取得，且题中点 \((x,y)\) 在内部，不能取到边界顶点的极值，所以
\[
-14<z<20
\]
故选 B.
\end{solution}

\begin{problem}
已知函数 \( f(x) = \begin{cases}
|\lg x|, & 0 < x \leqslant 10,\\
-\frac{1}{2}x + 6, & x > 10,
\end{cases} \) 若 \(a\)，\(b\)，\(c\) 互不相等，且 \( f(a) = f(b) = f(c) \)，则 \( abc \) 的取值范围是（\quad）

\choices
    {\((1,10)\)}
    {\((5,6)\)}
    {\((10,12)\)}
    {\((20,24)\)}
\end{problem}

\begin{answer}
C
\end{answer}

\begin{solution}
设 \(f(a)=f(b)=f(c)=r\). 当 \(0<x<1\) 时，\(f(x)=-\lg x\)；当 \(1\leqslant x\leqslant10\) 时，\(f(x)=\lg x\)；当 \(x>10\) 时，\(f(x)=-\dfrac12x+6\).

为了有三个互不相等的原像，必须取 \(0<r<1\)：此时三个原像分别来自上述三个区间，且为
\(x_1=10^{-r}\)，\(x_2=10^{r}\)，\(x_3=12-2r\).
当 \(r=0\) 或 \(r=1\) 时只有两个原像；当 \(r<0\) 或 \(r>1\) 时也不可能有三个原像. 因此
\[
abc=x_1x_2x_3=12-2r
\]
由 \(0<r<1\)，得 \(10<12-2r<12\)，所以 \(abc\in(10,12)\)，故选 C.
\end{solution}

\section{填空题}

\begin{problem}
圆心位于原点且与直线 \(x + y - 2 = 0\) 相切的圆的方程为\fillinblank{}.
\end{problem}

\begin{answer}
\(x^{2}+y^{2}=2\)
\end{answer}

\begin{solution}
圆心在原点，设半径为 \(r\). 圆与直线 \(x+y-2=0\) 相切，所以圆心到该直线的距离等于半径，即
\[
r=\frac{|0+0-2|}{\sqrt{1^{2}+1^{2}}}=\sqrt{2}
\]
因此圆的方程为
\[
x^{2}+y^{2}=r^{2}=2
\]
\end{solution}

\begin{problem}
设函数 \( y = f(x) \) 在区间 \([0,1]\) 上的图象是连续不断的一条曲线，且恒有 \( 0 \leqslant f(x) \leqslant 1 \)，可以用随机模拟方法近似计算由曲线 \( y = f(x) \) 及直线 \(x = 0\)，\(x = 1\)，\(y = 0\) 所围成部分的面积 \(S\)，先产生两组（每组 \(N\) 个）区间 \([0,1]\) 上的均匀随机数 \(x_1\)，\(x_2\)，\(\cdots\)，\(x_N\) 和 \(y_1\)，\(y_2\)，\(\cdots\)，\(y_N\)，由此得到 \(N\) 个点 \( (x_i, y_i) \)（\(i = 1\)，\(2\)，\(\cdots\)，\(N\)），再数出其中满足 \( y_i \leqslant f(x_i) \)（\(i = 1\)，\(2\)，\(\cdots\)，\(N\)）的点数 \(N_{1}\)，那么由随机模拟方法可得 \(S\) 的近似值为\fillinblank{}.
\end{problem}

\begin{answer}
\(\dfrac{N_{1}}{N}\)
\end{answer}

\begin{solution}
由于 \(0\leqslant f(x)\leqslant1\)，曲线及直线 \(x=0\)，\(x=1\)，\(y=0\) 围成的区域包含在单位正方形 \([0,1]\times[0,1]\) 内，其面积就是随机点落在区域内的概率乘以单位正方形面积. 对于均匀随机产生的 \(N\) 个点，落在该区域内的频率近似为
\[
\frac{N_{1}}{N}
\]
因为单位正方形面积为 \(1\)，所以区域面积 \(S\) 的近似值为
\[
S\approx\frac{N_{1}}{N}
\]
\end{solution}

\begin{problem}
一个几何体的正视图为一个三角形，则这个几何体可能是下列几何体中的\fillinblank{}. （填入所有可能的几何体前的编号）

\begin{circlelist}
\item 三棱锥；
\item 四棱锥；
\item 三棱柱；
\item 四棱柱；
\item 圆锥；
\item 圆柱.
\end{circlelist}
\end{problem}

\begin{answer}
\(\circled{1}\circled{2}\circled{3}\circled{5}\)
\end{answer}

\begin{solution}
三棱锥和四棱锥从适当方向观察时，底面或侧面投影可以形成三角形；三棱柱从垂直于其三角形底面的方向观察时，正视图可以是三角形；圆锥的轴截面为三角形，也可以得到三角形正视图. 四棱柱的正视图为四边形，圆柱的正视图为矩形，均不能为三角形. 因此可能的几何体编号为 \(1\)，\(2\)，\(3\)，\(5\)，即 \(\circled{1}\circled{2}\circled{3}\circled{5}\).
\end{solution}

\begin{problem}
在 \(\triangle ABC\) 中，\(D\) 为 \(BC\) 边上一点，\(BC = 3BD\)，\(AD = \sqrt{2}\)，\(\angle ADB = 135^{\circ}\). 若 \(AC = \sqrt{2}AB\)，则 \(BD =\fillinblank{}\).
\end{problem}

\begin{answer}
\(2+\sqrt{5}\)
\end{answer}

\begin{solution}
设 \(BD=x\)，则 \(BC=3x\)，从而 \(DC=BC-BD=2x\). 因为 \(B\)，\(D\)，\(C\) 三点共线，且 \(\angle ADB=135^{\circ}\)，所以 \(\angle ADC=45^{\circ}\).

在 \(\triangle ABD\) 中，由余弦定理，
\[
AB^{2}=AD^{2}+BD^{2}-2AD\cdot BD\cos135^{\circ}
=2+x^{2}+2x=x^{2}+2x+2
\]
在 \(\triangle ADC\) 中，
\[
AC^{2}=AD^{2}+DC^{2}-2AD\cdot DC\cos45^{\circ}
=2+4x^{2}-4x=4x^{2}-4x+2
\]
由 \(AC=\sqrt{2}AB\)，得
\[
4x^{2}-4x+2=AC^{2}=2AB^{2}=2x^{2}+4x+4
\]
整理得 \(x^{2}-4x-1=0\)，所以 \(x=2\pm\sqrt{5}\). 因为 \(x=BD>0\)，故 \(x=2+\sqrt{5}\)，即
\[
BD=2+\sqrt{5}
\]
\end{solution}

\section{解答题}

\begin{problem}
设等差数列 \(\{a_{n}\}\) 满足 \(a_{3} = 5\)，\(a_{10} = -9\).

\begin{enumerate}
\item 求 \(\{a_{n}\}\) 的通项公式；
\item 求 \(\{a_{n}\}\) 的前 \(n\) 项和 \(S_{n}\) 及使得 \(S_{n}\) 最大的序号 \(n\) 的值.
\end{enumerate}
\end{problem}

\begin{solution}
\begin{enumerate}
\item 设等差数列 \(\{a_n\}\) 的首项为 \(a_1\)，公差为 \(d\)，则 \(a_3=a_1+2d=5\)，\(a_{10}=a_1+9d=-9\). 两式相减得 \(7d=-14\)，所以 \(d=-2\)，进而 \(a_1=9\). 因此 \(a_n=a_1+(n-1)d=9-2(n-1)=11-2n\).

\item 前 \(n\) 项和为 \(S_n=\dfrac{n(a_1+a_n)}{2}=\dfrac{n(9+11-2n)}{2}=n(10-n)=25-(n-5)^2\). 由于 \(n\) 为正整数，当 \(n=5\) 时 \(S_n\) 取得最大值 \(25\). 所以使 \(S_n\) 最大的序号为 \(n=5\).
\end{enumerate}
\end{solution}

\begin{problem}
如图，已知四棱锥 \(P - ABCD\) 的底面为等腰梯形，\(AB \parallel CD\)，\(AC \perp BD\)，垂足为 \(H\)，\(PH\) 是四棱锥的高.

\begin{enumerate}
\item 证明：平面 \(PAC \perp\) 平面 \(PBD\)；
\item 若 \(AB = \sqrt{6}\)，\(\angle APB = \angle ADB = 60^{\circ}\)，求四棱锥 \(P - ABCD\) 的体积.
\end{enumerate}

\bitmapfigure[max width=0.34\linewidth]{img/2010/new_curriculum_liberal/q18_fig1.png}
\end{problem}

\begin{solution}
\begin{enumerate}
\item 因为 \(PH\perp\) 平面 \(ABCD\)，且 \(AC\) 在平面 \(ABCD\) 内，所以 \(AC\perp PH\). 又因为 \(AC\perp BD\)，且 \(PH\) 与 \(BD\) 是平面 \(PBD\) 内相交的两条直线，所以 \(AC\perp\) 平面 \(PBD\). 由于 \(AC\) 在平面 \(PAC\) 内，故平面 \(PAC\perp\) 平面 \(PBD\).

\item 设 \(AB=2u\)，\(CD=2v\)，底面等腰梯形的高为 \(h_0\). 建立直角坐标系，使
\(A(-u,0)\)，\(B(u,0)\)，\(D(-v,h_0)\)，\(C(v,h_0)\). 由 \(AC\perp BD\)，得
\((u+v,h_0)\cdot (-(u+v),h_0)=0\)，所以 \(h_0=u+v\).

由此可知对角线交点 \(H\) 的坐标为 \((0,u)\)，从而 \(AH=BH=\sqrt{2}u\). 在三角形 \(ADB\) 中，
\(\overrightarrow{DA}=(v-u,-u-v)\)，\(\overrightarrow{DB}=(v+u,-u-v)\)，所以
\[
\cos\angle ADB
=\frac{\overrightarrow{DA}\cdot\overrightarrow{DB}}{|DA||DB|}
=\frac{v}{\sqrt{u^2+v^2}}
\]
由 \(\angle ADB=60^\circ\)，得 \(v/\sqrt{u^2+v^2}=1/2\)，即 \(u^2=3v^2\). 又 \(AB=\sqrt{6}\)，故 \(u=\sqrt{6}/2\)，\(v=\sqrt{2}/2\)，于是 \(CD=\sqrt{2}\)，且底面高 \(h_0=(\sqrt{6}+\sqrt{2})/2\). 因此底面面积为
\[
S_{\text{底}}=\frac{AB+CD}{2}h_0
=\frac{(\sqrt{6}+\sqrt{2})^2}{4}=2+\sqrt{3}
\]

因为 \(AC\perp BD\)，所以 \(HA\perp HB\)，结合 \(AH=BH\) 及 \(PH\perp\) 平面 \(ABCD\)，可得 \(PA=PB\). 又 \(\angle APB=60^\circ\)，所以三角形 \(APB\) 为等边三角形，\(PA=AB=\sqrt{6}\). 在直角三角形 \(PHA\) 中，\(PH^2=PA^2-AH^2=6-3=3\)，故 \(PH=\sqrt{3}\). 四棱锥体积为
\[
V=\frac{1}{3}S_{\text{底}}PH
=\frac{1}{3}(2+\sqrt{3})\sqrt{3}
=1+\frac{2\sqrt{3}}{3}
\]
\end{enumerate}
\end{solution}

\begin{problem}
为调查某地区老人是否需要志愿者提供帮助，用简单随机抽样方法从该地区调查了 \(500\text{ 位}\)老年人，结果如下：

\begin{center}
\small
\setlength{\tabcolsep}{0.8em}
\renewcommand{\arraystretch}{1.2}
\begin{tabular}{|c|c|c|}
\hline
\makebox[3.6cm][l]{%
  \rule[-18.82pt]{0pt}{43.12pt}%
  \smash{\rlap{\raisebox{1.05cm}{\hspace{-0.3cm}\rotatebox{-19.75}{\rule{4.49cm}{0.4pt}}}}}%
  \rlap{\raisebox{0.74cm}{\hspace{2.6cm}性别}}%
  \rlap{\raisebox{-0.22cm}{\hspace{0.12cm}是否需要志愿者}}%
}
& 男 & 女 \\
\hline
需要 & \(40\) & \(30\) \\
\hline
不需要 & \(160\) & \(270\) \\
\hline
\end{tabular}
\end{center}

\begin{enumerate}
\item 估计该地区老年人中，需要志愿者提供帮助的老年人的比例；
\item 能否有 \(99\%\) 的把握认为该地区的老年人是否需要志愿者提供帮助与性别有关？
\item 根据（2）的结论，能否提供更好的调查方法来估计该地区老年人，需要志愿者帮助的老年人的比例？说明理由.
\end{enumerate}

附：
\begin{center}
\begin{tabular}{c c c c}
\(P(K^2\geqslant k)\) & \(0.050\) & \(0.010\) & \(0.001\) \\
\hline
\(k\) & \(3.841\) & \(6.635\) & \(10.828\)
\end{tabular}
\end{center}

\(\displaystyle K^{2} = \frac{n(ad - bc)^{2}}{(a + b)(c + d)(a + c)(b + d)}\)
\end{problem}

\begin{solution}
\begin{enumerate}
\item 样本中需要志愿者的老年人数为 \(40+30=70\)，所以可以估计该地区老年人中需要志愿者提供帮助的比例为
\(\dfrac{70}{500}=0.14\)，即 \(14\%\).

\item 令“是否需要志愿者提供帮助”与“性别”相互独立为原假设. 在给定的列联表中取 \(a=40\)，\(b=30\)，\(c=160\)，\(d=270\)，\(n=500\)，则
\[
K^2=\frac{500(40\times270-30\times160)^2}{(40+30)(160+270)(40+160)(30+270)}
\approx9.967
\]
由于 \(9.967>6.635\)，所以在犯错误的概率不超过 \(1\%\) 的前提下拒绝原假设. 因此有 \(99\%\) 的把握认为该地区老年人是否需要志愿者提供帮助与性别有关.

\item 可以按性别进行分层抽样，分别从男性老年人和女性老年人中随机抽取样本，并按两类老年人在该地区总体中的人数比例加权估计总体比例. 这是因为是否需要志愿者提供帮助与性别有关，分层后各层内的调查对象具有较强同质性，能够减少性别结构差异造成的抽样误差，从而得到更可靠的总体比例估计.
\end{enumerate}
\end{solution}

\begin{problem}
设 \(F_{1}\)，\(F_{2}\) 分别是椭圆 \(E: x^{2} + \frac{y^{2}}{b^{2}} = 1\)（\(0 < b < 1\)）的左，右焦点，过 \(F_{1}\) 的直线 \(l\) 与 \(E\) 相交于 \(A\)，\(B\) 两点，且 \(|AF_{2}|\)，\(|AB|\)，\(|BF_{2}|\) 成等差数列.

\begin{enumerate}
\item 求 \(|AB|\)；
\item 若直线 \(l\) 的斜率为 1，求 \(b\) 的值.
\end{enumerate}
\end{problem}

\begin{solution}
\begin{enumerate}
\item 设椭圆的半焦距为 \(c=\sqrt{1-b^2}\)，则 \(F_1(-c,0)\)，\(F_2(c,0)\). 直线 \(l\) 经过椭圆内的焦点 \(F_1\)，所以 \(F_1\) 在线段 \(AB\) 上. 由椭圆的定义，\(|AF_1|+|AF_2|=2\)，\(|BF_1|+|BF_2|=2\). 因而
\(|AB|=|AF_1|+|F_1B|=4-|AF_2|-|BF_2|\).

由于 \(|AF_2|\)，\(|AB|\)，\(|BF_2|\) 成等差数列，有 \(2|AB|=|AF_2|+|BF_2|\). 联立得 \(3|AB|=4\)，所以 \(|AB|=\dfrac{4}{3}\).

\item 当直线 \(l\) 的斜率为 \(1\) 时，其方程为 \(y=x+c\). 将它代入椭圆方程，得
\[
(1+b^2)x^2+2cx+c^2-b^2=0
\]
设两个交点的横坐标为 \(x_1\)，\(x_2\). 因为 \(c^2=1-b^2\)，该方程的判别式为
\[
\Delta=4c^2-4(1+b^2)(c^2-b^2)=8b^4
\]
所以 \(|x_1-x_2|=\dfrac{\sqrt{\Delta}}{1+b^2}=\dfrac{2\sqrt{2}b^2}{1+b^2}\). 两交点在斜率为 \(1\) 的直线上，故
\[
|AB|=\sqrt{2}|x_1-x_2|=\frac{4b^2}{1+b^2}
\]
由第一问，\(\dfrac{4b^2}{1+b^2}=\dfrac{4}{3}\)，解得 \(b^2=\dfrac12\). 因为 \(b>0\)，所以 \(b=\dfrac{\sqrt{2}}{2}\).
\end{enumerate}
\end{solution}

\begin{problem}
设函数 \( f(x) = x(\e^x - 1) - ax^2 \).

\begin{enumerate}
\item 若 \(a = \frac{1}{2}\)，求 \(f(x)\) 的单调区间；
\item 若当 \(x \geqslant 0\) 时，\(f(x) \geqslant 0\)，求 \(a\) 的取值范围.
\end{enumerate}
\end{problem}

\begin{solution}
\begin{enumerate}
\item 当 \(a=\dfrac12\) 时，\(f'(x)=\e^x(x+1)-1-x=(x+1)(\e^x-1)\). 当 \(x<-1\) 时，\(x+1<0\) 且 \(\e^x-1<0\)，所以 \(f'(x)>0\)；当 \(-1<x<0\) 时，\(x+1>0\) 且 \(\e^x-1<0\)，所以 \(f'(x)<0\)；当 \(x>0\) 时，\(f'(x)>0\). 因此 \(f(x)\) 在 \((-\infty,-1)\) 和 \((0,+\infty)\) 上单调递增，在 \((-1,0)\) 上单调递减.

\item 当 \(x>0\) 时，\(f(x)\geqslant0\) 等价于
\(\e^x-1-ax\geqslant0\)，即 \(a\leqslant\dfrac{\e^x-1}{x}\). 由 \(\e^x\geqslant1+x\)，得 \(\dfrac{\e^x-1}{x}\geqslant1\)，所以当 \(a\leqslant1\) 时，对任意 \(x\geqslant0\) 都有
\[
f(x)=x(\e^x-1-ax)\geqslant x(\e^x-1-x)\geqslant0
\]

反之，若 \(a>1\)，由于 \(\displaystyle\lim_{x\to0^+}\dfrac{\e^x-1}{x}=1\)，可取充分小的正数 \(x\)，使 \(\dfrac{\e^x-1}{x}<a\)，此时 \(f(x)<0\). 因此 \(a\) 的取值范围为 \(a\leqslant1\).
\end{enumerate}
\end{solution}

\section{选考题：解答题}

\begin{problem}
如图，已知圆上的弧 \(\widehat{AC} = \widehat{BD}\)，过 \(C\) 点的圆的切线与 \(BA\) 的延长线交于 \(E\) 点，证明：

\begin{enumerate}
\item \(\angle ACE = \angle BCD\)；
\item \(BC^{2} = BE\cdot CD\).
\end{enumerate}

\bitmapfigure[max width=0.65\linewidth]{img/2010/new_curriculum_liberal/q22_fig1.png}
\end{problem}

\begin{solution}
\begin{enumerate}
\item 由切线与弦所夹的角等于弦所对的圆周角，得 \(\angle ACE=\angle ABC\). 又因为弧 \(\widehat{AC}\) 与弧 \(\widehat{BD}\) 相等，所以它们所对的圆周角相等，即 \(\angle ABC=\angle BCD\). 因此 \(\angle ACE=\angle BCD\).

\item 因为 \(E\)，\(A\)，\(B\) 三点共线，所以 \(\angle EBC=\angle ABC=\angle BCD\). 又由切线与弦所夹的角定理，\(\angle ECB=\angle CDB\). 所以 \(\triangle EBC\sim\triangle BCD\)，对应关系为 \(E\leftrightarrow B\)，\(B\leftrightarrow C\)，\(C\leftrightarrow D\). 从而 \(\dfrac{BE}{BC}=\dfrac{BC}{CD}\)，即 \(BC^2=BE\cdot CD\).
\end{enumerate}
\end{solution}

\begin{problem}
已知直线 \(C_{1}: \begin{cases}
x = 1 + t\cos \alpha,\\
y = t\sin \alpha,
\end{cases}\) （\(t\) 为参数），圆 \(C_{2}: \begin{cases}
x = \cos \theta,\\
y = \sin \theta.
\end{cases}\) （\(\theta\) 为参数）.

\begin{enumerate}
\item 当 \(\alpha = \frac{\pi}{3}\) 时，求 \(C_{1}\) 与 \(C_{2}\) 的交点坐标；
\item 过坐标原点 \(O\) 作 \(C_{1}\) 的垂线，垂足为 \(A\)，\(P\) 为 \(OA\) 的中点，当 \(\alpha\) 变化时，求点 \(P\) 轨迹的参数方程，并指出它是什么曲线.
\end{enumerate}
\end{problem}

\begin{solution}
\begin{enumerate}
\item 当 \(\alpha=\dfrac{\pi}{3}\) 时，直线 \(C_1\) 的参数方程为 \(x=1+\dfrac{t}{2}\)，\(y=\dfrac{\sqrt{3}}{2}t\). 圆 \(C_2\) 的方程为 \(x^2+y^2=1\). 代入得
\[
\left(1+\frac{t}{2}\right)^2+\frac{3t^2}{4}=1
\]
即 \(t^2+t=0\)，所以 \(t=0\) 或 \(t=-1\). 当 \(t=0\) 时交点为 \((1,0)\)；当 \(t=-1\) 时交点为 \(\left(\dfrac12,-\dfrac{\sqrt{3}}2\right)\).

\item 直线 \(C_1\) 经过点 \((1,0)\)，方向向量为 \((\cos\alpha,\sin\alpha)\). 从原点 \(O\) 到直线 \(C_1\) 的垂足 \(A\) 对应的参数为
\(t=-\cos\alpha\)，因此
\[
A=(1,0)-\cos\alpha(\cos\alpha,\sin\alpha)
=\left(\sin^2\alpha,-\sin\alpha\cos\alpha\right)
\]
由于 \(P\) 是 \(OA\) 的中点，点 \(P\) 的参数方程为
\[
\begin{cases}
x=\dfrac12\sin^2\alpha,\\
y=-\dfrac12\sin\alpha\cos\alpha.
\end{cases}
\]
利用二倍角公式，得 \(x-\dfrac14=-\dfrac14\cos2\alpha\)，\(y=-\dfrac14\sin2\alpha\)，所以
\[
\left(x-\frac14\right)^2+y^2=\frac1{16}
\]
故点 \(P\) 的轨迹是圆心为 \(\left(\dfrac14,0\right)\)、半径为 \(\dfrac14\) 的圆.
\end{enumerate}
\end{solution}

\begin{problem}
设函数 \( f(x) = |2x - 4| + 1 \).

\begin{enumerate}
\item 画出函数 \( y = f(x) \) 的图象；
\item 若不等式 \( f(x) \leqslant ax \) 的解集非空，求 \( a \) 的取值范围.
\end{enumerate}

\bitmapfigure[max width=0.32\linewidth]{img/2010/new_curriculum_liberal/q24_fig1.png}
\end{problem}

\begin{solution}
\begin{enumerate}
\item 函数可写成
\[
f(x)=\begin{cases}
-2x+5,&x<2,\\
2x-3,&x\geqslant2.
\end{cases}
\]
所以函数图象是开口向上的折线，顶点为 \((2,1)\)，左侧斜率为 \(-2\)，右侧斜率为 \(2\). 例如图象经过 \((0,5)\) 和 \((\frac52,2)\)，据此在给定坐标系中作出两条射线即可.

\item \(x=0\) 时，\(f(0)=5\)，不满足不等式. 若 \(x>0\)，则
\[
\frac{f(x)}{x}=
\begin{cases}
\dfrac{5}{x}-2,&0<x<2,\\
2-\dfrac{3}{x},&x\geqslant2.
\end{cases}
\]
当 \(0<x<2\) 时，\(\dfrac{5}{x}-2>\dfrac12\)；当 \(x\geqslant2\) 时，\(2-\dfrac{3}{x}\geqslant\dfrac12\). 因而若不等式在正数 \(x\) 处有解，则 \(a\geqslant\dfrac12\). 反之，当 \(a\geqslant\dfrac12\) 时取 \(x=2\)，有 \(f(2)=1\leqslant2a\)，所以必有解.

若 \(x<0\)，则 \(f(x)=-2x+5\)，不等式除以负数 \(x\) 后变为 \(a\leqslant-2+\dfrac{5}{x}<-2\). 反之，当 \(a<-2\) 时取 \(x=\dfrac{5}{a+2}<0\)，可得 \(a=-2+\dfrac{5}{x}\)，原不等式成立. 所以 \(a\) 的取值范围为
\(a<-2\) 或 \(a\geqslant\dfrac12\)，即 \((-\infty,-2)\cup[\dfrac12,+\infty)\).
\end{enumerate}
\end{solution}
